% Bilaga F backs a design choice that only the teacher's notes make (a
% live lecture with questions instead of the flipped classroom), so
% sokprotokoll.tex inputs it only when the notes are shown.
%
% ======================================================================
% Bilaga F: the flipped-classroom review of the course-evaluation report
% (research/evaluation/delutvärdering-2026/bilaga-flippat.tex, d536059),
% adapted to the deck's appendix form.  Only the literature review is
% used; nothing from the students' survey answers.
%
% Author's note (verification and reproduction):
%   scholar session : flipped-classroom-learning (not rerun here)
%   exports         : research/evaluation/delutvärdering-2026/
%                     litteratursokning/hits.csv (classified, 845 unique
%                     records) and hits-table.tex (bearing-only longtable,
%                     copied as litteratursokning/recap-flippat-full.tex)
%   searched        : 2026-10-07; providers openalex, wos, scopus, ieee
%                     (scopus standard view, no abstracts);
%                     s2: key rejected (HTTP 403); dblp: every request
%                     answered with an Anubis bot-challenge page (non-JSON)
%   queries         : S1 S2 S3 C1 C2 C3 P1 P2 + P1w (P1 with truncation, on
%                     ieee/wos/scopus); verbatim in the header of the
%                     report's litteratur.bib, copied into
%                     ltnotes-repetition.bib.  In the text C1-C3 are M1-M3.
%                     The first round (same queries) was rerun: its WoS
%                     cells used TS=(FC) AND ... so TS covered only the
%                     first group (400 errors), and OpenAlex rejected
%                     wildcards inside phrases.  Coverage of the redone
%                     search checked by normalised title against the first
%                     round: all first-round items present after adding
%                     P1w (2026-10-07).
%   limit           : -l 50 per query and provider
%   classification  : scholar llm context + llm classify --no-examples,
%                     model openai-codex/gpt-5.6-sol, tags supports-claim /
%                     refutes-claim / qualifies-claim / adjacent-subtopic /
%                     off-topic-false-hit; every refutes-claim row and every
%                     cited source read; one row below confidence 0.6.
%                     Cited sources recorded with scholar sessions decide
%                     --keep and theme tags metaanalys, metodkritik,
%                     programmering, preferens.
%   full texts      : scholar PDF cache (Låg & Sæle via ERIC EJ1229661,
%                     van Alten via Utrecht repository, Hew et al. 2021
%                     via figshare, Kapur, Almassri) and Europe PMC XML
%                     (Deslauriers PMC6765278, Hew & Lo PMC5855972, Jensen
%                     PMC4353080, Boedeker PMC11933506).
%   provenance      : ltnotes-repetition.bib, the 20 entries after the
%                     "Bilaga F" header, each block copied verbatim with a
%                     FOUND-VIA (here) line; quotes of LagSaele2019 and
%                     vanAlten2019 spot-checked 2026-10-10.
%   changes         : section structure -> chapter; \parencite -> \autocite;
%                     references to the report's section 3.6 and the
%                     course-specific discussion rewritten for the recap
%                     deck; the report's sentence on the students' wish
%                     for lectures (it rests on the survey) removed.

\chapter{Ger ett flippat klassrum bättre lärande än föreläsningar?}
\label{app:recap-flippat}
\chapterprecis{Författaren har ännu inte granskat resultaten i den här
  bilagan i sin helhet.}

I genomsnitt ja, men lite: ett flippat klassrum ger omkring en femtedels
standardavvikelse bättre lärande än föreläsningsbaserad undervisning när
effekten korrigeras för publiceringsbias, effekten varierar mycket
mellan kurser, och den krymper eller försvinner när även den
föreläsningsbaserade undervisningen har aktiva moment.

Kursen var tidigare upplagd som ett flippat klassrum: i ett flippat
klassrum tar studenterna del av genomgången, oftast som video, före
passet, och passet används till att arbeta med stoffet.  Sedan 2026
bygger kursen i stället på en föreläsning live, med frågor som alla
svarar på under föreläsningen, följd av övning och laboration, och den
här föreläsningen är byggd just så, med en fråga först i varje avsnitt
(den första i \cref{ovn:recap-bankomat}).
Bytet vilar på ett påstående: att ett flippat klassrum i genomsnitt ger
något bättre studieresultat än föreläsningsbaserad
undervisning\autocite{LagSaele2019,vanAlten2019}, men att effekten är
liten och att den krymper när även den traditionella undervisningen har
aktiva moment.  Bilagan prövar påståendet genom att ställa frågan om
själva sakförhållandet: \emph{ger ett flippat klassrum bättre lärande än
föreläsningsbaserad undervisning, i högre utbildning i allmänhet och i
inledande programmering i synnerhet?}  Till den hör tre följdfrågor: hur
stor och hur säker effekten är, vilka drag i upplägget som bär den, och
hur studenternas nöjdhet och upplevda lärande förhåller sig till det
faktiska lärandet.  Påståendet som prövas är \enquote{i genomsnitt något
bättre}, inte \enquote{alltid bättre}.  Frågorna under föreläsningen har
sitt eget underlag i \cref{app:recap-retrieval,app:recap-polls}.

\section{Metod}
\label{sec:recap-flippat-metod}

Sökningen gjordes den 7 oktober 2026 med verktyget \texttt{scholar}, som
skickar samma fråga till flera bibliografiska databaser och sparar frågor
och träffar.  Fyra databaser svarade: OpenAlex, Web of Science, Scopus och
IEEE Xplore; Scopus lämnade bara grunduppgifter, utan sammanfattningar.
Två databaser föll bort: DBLP svarade på varje fråga med en sida som
kontrollerar att besökaren inte är en robot, och Semantic Scholar avvisade
verktygets nyckel.  Deras kolumner saknas därför i
\cref{tab:recap-flippat-fragor}; de betyder inte noll träffar.

Frågorna byggs av samma begreppsblock.  Flippat klassrum (FC) är
\enquote{flipped classroom} OR \enquote{inverted classroom} OR
\enquote{flipped learning}; utfallen är achievement, \enquote{learning
outcomes}, performance och grades; översikterna är meta-analysis,
\enquote{systematic review} och \enquote{literature review}; ämnet är
programming, \enquote{computer science} och computing.  Stödfrågorna (S)
söker översikter och jämförelser.  Motfrågorna (M) återanvänder blocken och
lägger till det som skulle försvaga påståendet: publiceringsbias,
heterogenitet, moderatorer, begränsningar, \enquote{no significant
difference}, negativa resultat och motstånd.  Två frågor (P) gäller
studenternas preferenser och upplevda lärande.  Web of Science fick varje
fråga som \texttt{TS=(\dots)}, Scopus som \texttt{TITLE-ABS-KEY(\dots)},
OpenAlex och IEEE Xplore som boolesk fritext.  Fråga P1 kördes med
trunkering (\texttt{lecture*}, \texttt{preference*}) i de databaser som
tillåter det och med utskrivna böjningsformer i OpenAlex.  Varje fråga gav
högst 50 träffar per databas, i databasens relevansordning.

\begin{table}[htbp]
  \centering\footnotesize
  \caption{Sökfrågor och antal träffar per databas, 7 oktober 2026, med
    verktyget \texttt{scholar}. S = stödfråga, M = motfråga, P = fråga om
    preferens och upplevt lärande. FC står för (\enquote{flipped
    classroom} OR \enquote{inverted classroom} OR \enquote{flipped
    learning}); (s) betyder att både singular och plural söktes.
    Databaser: OpenAlex (OA; inte \foreignlanguage{english}{open access}),
    Web of Science (WoS), Scopus och IEEE Xplore (IEEE). Högst 50 träffar
    per databas och fråga, så 50 betyder \enquote{minst 50}. DBLP och
    Semantic Scholar svarade inte och saknas.}
  \label{tab:recap-flippat-fragor}
  \begin{tabular}{@{}l>{\raggedright\arraybackslash}p{0.62\linewidth}rrrr@{}}
    \toprule
    & Fråga & OA & WoS & Scopus & IEEE \\
    \midrule
    S1 & FC AND (meta-analysis OR \enquote{systematic review}) AND
      (achievement OR \enquote{learning outcomes} OR performance)
      & 50 & 50 & 50 & 18 \\
    S2 & FC AND (programming OR \enquote{computer science} OR computing)
      AND (achievement OR \enquote{learning outcomes} OR performance OR
      grades) & 50 & 50 & 50 & 50 \\
    S3 & FC AND (programming OR \enquote{computer science} OR computing)
      AND (meta-analysis OR \enquote{systematic review} OR
      \enquote{literature review}) & 50 & 50 & 47 & 28 \\
    M1 & FC AND (meta-analysis OR \enquote{systematic review}) AND
      (\enquote{publication bias} OR heterogeneity OR \enquote{no
      significant difference} OR moderator(s) OR limitations)
      & 50 & 50 & 50 & 6 \\
    M2 & FC AND (programming OR \enquote{computer science} OR computing)
      AND (\enquote{no significant difference} OR \enquote{no difference}
      OR challenges OR resistance OR negative) & 50 & 50 & 50 & 50 \\
    M3 & S3 AND (\enquote{no significant difference} OR heterogeneity OR
      limitations OR \enquote{mixed results}) & 50 & 8 & 1 & 1 \\
    P1 & (FC OR \enquote{active learning}) AND lecture(s) AND
      (\enquote{student perception(s)} OR satisfaction OR preference(s) OR
      resistance OR \enquote{feeling of learning}) & 50 & 50 & 50 & 50 \\
    P2 & (FC OR \enquote{active learning}) AND (\enquote{perceived
      learning} OR \enquote{feeling of learning} OR \enquote{self-reported
      learning}) AND (lecture(s) OR passive) & 50 & 33 & 40 & 6 \\
    \bottomrule
  \end{tabular}
\end{table}

Sökningen gav 845 unika träffar.  En språkmodell klassade var och en mot en
beskrivning av frågan, som stöder påståendet, nyanserar det, motsäger det,
angränsande ämne eller felträff, med en angiven säkerhet; en träff kunde
få både \enquote{stöder} och \enquote{nyanserar}.  Varje träff som modellen
klassade som motsägande lästes på sammanfattningsnivå, liksom varje källa
som bilagan bygger på; bara en träff hade en säkerhet under 0,6.

Källorna valdes i fyra lager: först de metaanalyser som jämför flippat
klassrum med föreläsningsbaserad undervisning över flera ämnen, sedan
granskningar av metaanalysernas metod och kontrollerade jämförelser som
talar emot eller nyanserar, sedan det som rör programmering, datavetenskap
eller matematik, och sist studier av preferens och upplevt lärande.  Varje
källa lästes i fulltext där en öppen version fanns, hos förlaget, i ett
lärosätes arkiv eller i Europe PMC; annars lästes sammanfattningen.
\Cref{tab:recap-flippat-kallor} anger vilket för varje källa.  Siffror ur
en fulltext anges med sidnummer; fyra fulltexter fanns bara utan
sidnumrering, och för dem anges sidan bara där den framgår.  Siffror som
bara kommer från en sammanfattning sägs komma därifrån.

\section{Resultat}
\label{sec:recap-flippat-resultat}

\subsection{Effekten i högre utbildning}

Den bredaste metaanalysen, av \textcite{LagSaele2019}, jämför flippat
klassrum med traditionell föreläsningsbaserad undervisning i 272 studier
från alla nivåer och ämnen.  Den sammanvägda effekten på lärande är
\(g = 0{,}35\), med 95\,\% konfidensintervall 0,31--0,40
\autocite[1]{LagSaele2019}.  Med bara de studier som har tillräcklig
statistisk styrka blir den 0,24, och korrigerad för saknade små studier
med metoden \foreignlanguage{english}{trim and fill} 0,17; författarna
sammanfattar: \enquote{the true weighted mean effect is somewhere around
one fifth of a standard deviation}\autocite[9]{LagSaele2019}.  Andelen
godkända är 77,6\,\% mot 75,9\,\%\autocite[7]{LagSaele2019}.
\textcite[s.~8 och tabell~3, s.~10]{vanAlten2019} finner \(g = 0{,}36\)
(0,28--0,44) över 114 studier i gymnasium och högre utbildning.

Övriga metaanalyser pekar åt samma håll.  Enligt sammanfattningarna finner
\textcite{Cheng2019} \(g = 0{,}19\) över 55 publikationer,
\textcite{Strelan2020} \(g = 0{,}50\) över 198 studier, och
\textcite{Bredow2021}, med 317 studier i högre utbildning, signifikanta
fördelar för flippat klassrum i sju av åtta utfall, med \(g\) mellan 0,20
och 0,53.  Inom vårdutbildningar finner \textcite{HewLo2018} en
standardiserad medelvärdesskillnad på 0,33 (0,21--0,46).  En metaanalys av
15 metaanalyser i högre utbildning ger \(g = 0{,}45\) (0,37--0,53), och
0,37 efter korrigering för publiceringsbias\autocite[140--141]{Hew2021second};
dess tabell över de ingående analyserna bekräftar Chengs och Strelans
siffror för högre utbildning, 0,193 och 0,48\autocite[136]{Hew2021second}.
En liten metaanalys inom farmaci, med fem studier, finner däremot ingen
signifikant skillnad i tentamensresultat och talar om \enquote{minimal
gains}\autocite{Gillette2018}.

\subsection{Spridning, publiceringsbias och metodkvalitet}

Spridningen mellan studierna är mycket stor.  Andelen av variationen som
inte beror på slumpen, \(I^2\), är 83\,\% hos
\textcite[7]{LagSaele2019}, 88\,\% hos \textcite[tabell~3,
s.~10]{vanAlten2019} och 91\,\% i den metaanalys som
\textcite[7]{Kapur2022} gör.  Kapur m.fl.\ samlar också 46 metaanalyser
vars sammanvägda effekter ligger mellan 0,19 och
2,29\autocite[2]{Kapur2022}.  Publiceringsbias är belagd: små studier drar
upp effekten hos \textcite[9]{LagSaele2019}, och i metaanalysen av
metaanalyser ger både rangkorrelation och Eggers test tecken på
bias\autocite[141]{Hew2021second}.  En granskning av 19 metaanalyser
finner enligt sammanfattningen att bland annat sökning i grå litteratur,
bedömning av risk för bias, kontroll för studenternas förkunskaper och för
att samma lärare undervisar båda grupperna \enquote{were inadequately
conducted or reported}\autocite{Hew2021methodology}.

\subsection{Vad effekten beror på}

Två drag i upplägget har stöd.  Quiz på förberedelserna höjer effekten med
\(g = 0{,}19\), och en kortad lektionstid sänker den med
0,26\autocite[9]{vanAlten2019}; \textcite{HewLo2018} finner också större
effekt när passet inleds med quiz.  Hos \textcite[7]{LagSaele2019} är
skillnaden för prov på förberedelserna bara marginell (0,40 mot 0,31,
\(p = 0{,}051\)), och sociala aktiviteter i passet gör ingen påvisbar
skillnad.

Om det är det aktiva arbetet som bär effekten är omstritt.  Enligt
sammanfattningen ser \textcite{Strelan2020} möjligheten till strukturerat
aktivt lärande som den främsta orsaken.  \textcite{Kapur2022} kodar i
stället vad som faktiskt gjordes i 173 studier och finner att aktivt
lärande ofta saknades i de flippade passen och inte tillförde något när
det fanns; fördelen tycks komma av mer exponering och tid med stoffet.
De skriver: \enquote{As more active learning was incorporated in
traditional instruction, the effect of flipped learning over traditional
instruction tended to reduce, and in several cases, even
reverse}\autocite[12]{Kapur2022}.  Deras redovisning har dock en
motsägelse: på s.~10 kallas skillnaden mellan pass med och utan aktivt
lärande (\(g = 0{,}42\) mot 0,33) obefintlig, samtidigt som testet anges
till \(t = 4{,}11\), \(p < 0{,}001\).  \textcite{Jensen2015} jämförde i
ett kvasiexperiment två grupper i en biologikurs, 53 och 55 studenter med
samma lärare, där båda arbetade aktivt och bara ordningen skilde: lärandet
och nöjdheten blev desamma, och författarna drar slutsatsen att vinsterna
\enquote{are most likely a result of the active-learning style of
instruction rather than the order in which the instructor participated in
the learning process}.

\subsection{Programmering och datavetenskap}

Underlaget för programmering är tunt men pekar åt samma håll.  Enligt
sammanfattningen är effekten för IT-ämnen den svagaste hos
\textcite{Strelan2020}, \(g = 0{,}30\) över 14 studier.  Hos
\textcite[tabell~4, s.~11]{vanAlten2019} skiljer sig de formella
vetenskaperna, som artikeln inte närmare definierar, inte från
naturvetenskaperna (\(-0{,}035\), \(p = 0{,}698\)), och hos
\textcite[7]{LagSaele2019} är effekten för naturvetenskap, teknik och
matematik 0,32.  Den enda metaanalysen för programmering som sökningen
fann, av \textcite[267, 276]{AlmassriZaharudin2023}, väger samman 14
studier till \(g = 0{,}56\) med \(I^2 = 90{,}1\,\%\).  Den är svag:
konfidensintervallet anges olika i sammanfattningen och i texten
(0,33--0,79 mot 0,32--0,81, s.~275), och publiceringsbias bedöms bara med
trattdiagram och felsäkert \(N\).  Två enskilda studier finner, enligt
sammanfattningarna, ingen fördel eller en nackdel: inga signifikanta
skillnader i inledande programmering över flera år\autocite{Tyler2018},
och mindre lärandeökning med flippat klassrum i statistisk programmering,
som författarna förklarar med missuppfattningar som befästs under
självstudierna\autocite{Liu2026}.  En tidig översikt av 32 artiklar om
flippat klassrum i datavetenskap efterlyser kontrollerade
experiment\autocite{Giannakos2014}, och en översikt för matematik finner
stöd men inkonsekventa resultat, beroende på upplägg och
studentgrupp\autocite{Jones2026}.

\subsection{Nöjdhet, preferens och upplevt lärande}

I genomsnitt påverkar upplägget inte nöjdheten.  Hos
\textcite[9]{LagSaele2019} är effekten på nöjdhet 0,16 och efter
korrigering för saknade studier \(-0{,}03\), och hos
\textcite[tabell~3, s.~10]{vanAlten2019} 0,05, inte signifikant.  Låg och
Sæle påpekar att flera studier har effekter till föreläsningarnas fördel
och att \enquote{students may dislike an educational intervention even
though it improves learning}\autocite[12]{LagSaele2019}.  I
vårdutbildningarna föredrog däremot fler studenter det flippade
klassrummet\autocite{HewLo2018}.

\textcite{Deslauriers2019} lottade studenter i en inledande fysikkurs
mellan aktivt undervisade lektioner och flytande föreläsningar med samma
material; det aktiva var arbete under passet, inte förberedelser före,
så upplägget är inte flippat.  På provet låg de aktivt undervisade
studenterna nästan en halv standardavvikelse högre (0,46), men de
upplevde att de lärt sig mer än en halv standardavvikelse mindre (0,56)
och föredrog föreläsningarna: \enquote{students in the active classroom
learn more, but they feel like they learn
less}\autocite[19251]{Deslauriers2019}.  \textcite{Boedeker2025} gjorde
ett liknande försök med 146 läkarstudenter och fann ett högre
provresultat (0,27 standardavvikelser) men också ett högre upplevt
lärande (0,56) i de aktiva passen; de förklarar skillnaden mot
Deslauriers m.fl.\ med studenternas stadium eller innehållet.  Två
studier inom datavetenskap mäter bara upplevelsen, enligt
sammanfattningarna: i en flippad IT-kurs var studenterna nöjda med
upplägget vid kursens slut\autocite{Elliott2014}, och i inledande
programmering uppskattade studenterna måttligt aktiva tekniker som
livekodning mer än mycket aktiva som kodning under
passet\autocite{SoosaiRaj2018}.

\section{Diskussion och begränsningar}
\label{sec:recap-flippat-diskussion}

Det som talar för påståendet är att alla breda metaanalyser finner en
positiv effekt, oberoende av varandra och med delvis olika studier.  Det
som talar emot ett starkare påstående är att effekten är liten när den
korrigeras för publiceringsbias, att spridningen mellan studierna är så
stor att genomsnittet säger lite om en enskild kurs, att metaanalyserna
har kända metodbrister, och att en kontrollerad jämförelse och en
metaanalys som kodar vad som faktiskt gjordes i
passen\autocite{Jensen2015,Kapur2022} pekar på att fördelen krymper eller
försvinner när den icke-flippade undervisningen också är aktiv.  Vad som
bär effekten är omstritt; quiz på förberedelserna och bibehållen
lärarledd tid är de drag som har stöd i mer än en studie.

För den här föreläsningen betyder det att jämförelsen i de flesta
studierna, flippat klassrum mot föreläsning med hemuppgifter, inte är den
som görs här.  Kursens upplägg är föreläsning, live och inspelad, med
frågor under föreläsningen, följd av övning och laboration: en
traditionell ordning med aktiva moment under och efter genomgången.  Mot
ett sådant upplägg talar underlaget för en liten eller ingen skillnad,
och de drag som verkar bära effekten, uppföljda förberedelser och
lärarledd tid, kan finnas i båda; frågorna under föreläsningen är av samma
slag som de quiz som höjer effekten i de flippade klassrummen, och deras
eget underlag finns i \cref{app:recap-retrieval}.  Att studenter föredrar
eller inte föredrar ett upplägg säger lite om lärandet: nöjdheten beror i
genomsnitt inte på upplägget, och upplevt och faktiskt lärande kan gå
isär, även om de inte alltid gör det.

Underlaget har begränsningar.  Två av sex databaser föll bort, och ACM:s
konferenser om datavetenskaplig undervisning nås bara i den mån de andra
databaserna indexerar dem.  Varje fråga gav högst 50 träffar per databas,
så källor längre ned i relevansordningen kan ha missats.  Klassningen
gjordes av en språkmodell; bara de motsägande träffarna och de citerade
källorna lästes.  Elva av de tjugo källorna lästes bara som sammanfattning
(\cref{tab:recap-flippat-kallor}), bland dem fyra av metaanalyserna; för
Deslauriers m.fl.\ kunde bara en sida beläggas.  De flesta studierna är
inte randomiserade, och många gäller en kurs första försök med flippat
klassrum\autocite[12]{LagSaele2019}.  Tre metaanalyser av metaanalyser
från 2025 och 2026 hittades men kunde inte läsas.

\section{Slutsats}
\label{sec:recap-flippat-slutsats}

Ja, i genomsnitt ger ett flippat klassrum något bättre lärande än
föreläsningsbaserad undervisning, men effekten är liten, omkring en
femtedel av en standardavvikelse när den korrigeras för publiceringsbias,
och den varierar så mycket mellan kurser att den inte kan tas för given i
en enskild kurs.  För programmering pekar det tunna underlaget åt samma
håll.  Effekten tycks bäras av drag som uppföljda förberedelser och
bibehållen lärarledd tid snarare än av själva ordningen; om aktivt arbete
i passet bidrar är omstritt.  När den traditionella undervisningen har
aktiva moment är skillnaden liten eller ingen.  Nöjdheten påverkas i
genomsnitt inte; att studenter föredrar föreläsningar trots att de lär sig
mer av aktiv undervisning är visat i ett randomiserat försök men inte i
ett annat.  Att den här föreläsningen är en genomgång live med frågor, och
inte ett flippat klassrum, har alltså stöd i den meningen att underlaget
inte visar att ett flippat klassrum vore bättre än en föreläsning med
aktiva moment; att ett flippat klassrum är \enquote{bättre för
lärandet} utan förbehåll håller däremot inte.

\section{Fortsatt arbete}
\label{sec:recap-flippat-fortsatt}

Den här sökningen lämnade följande ogjort.  DBLP och Semantic Scholar bör
frågas när de svarar, och ACM:s digitala bibliotek bör sökas direkt.
Fulltexterna till \textcite{Strelan2020}, \textcite{Bredow2021},
\textcite{Cheng2019}, \textcite{Hew2021methodology}, \textcite{Liu2026},
\textcite{Giannakos2014}, \textcite{Gillette2018}, \textcite{Tyler2018},
\textcite{Jones2026}, \textcite{Elliott2014} och \textcite{SoosaiRaj2018}
behöver läsas, i första hand Strelan m.fl.\ (siffran för IT) och Bredow
m.fl.\ (vilket utfall som inte var signifikant).  Tre metaanalyser av
metaanalyser från 2025 och 2026 hittades men har inte lästs alls, inte
ens som sammanfattning: Ma m.fl.\ (2025),\footnote{Ma, Fan och Ning,
\enquote{The effectiveness of the flipped classroom: A second-order
meta-analysis}, \emph{International Journal of Educational Research},
2025, \url{https://doi.org/10.1016/j.ijer.2025.102855}.} Uzun m.fl.\
(2026)\footnote{Uzun, Kaya, Ünal och Mutlu-Bayraktar, \enquote{Impact of
flipped classrooms on higher education student outcomes: A second-order
meta-analysis}, \emph{Studies in Educational Evaluation}, 2026,
\url{https://doi.org/10.1016/j.stueduc.2026.101653}.} och Wang m.fl.\
(2026).\footnote{Wang, Wang, Mo och Shen, \enquote{Flipped classroom and
learning outcomes: a second-order meta-analysis}, \emph{Educational
Technology Research and Development}, 2026,
\url{https://doi.org/10.1007/s11423-026-10686-z}.} De kan ändra
storleken på effekten men knappast riktningen.

Sökningen fann inga randomiserade jämförelser i inledande programmering,
och i synnerhet inga mellan ett flippat klassrum och en icke-flippad kurs
som också är aktiv, som den Jensen m.fl.\ gjorde i biologi.  Den fann
inte heller någon studie som mäter både faktiskt och upplevt lärande hos
programmeringsnybörjare.  Om en sådan studie inte fann någon skillnad
mellan flippat klassrum och föreläsning med frågor och övningar, skulle
kursen kunna säga att inget av uppläggen är bättre för lärandet; fann den
en tydlig fördel för det flippade klassrummet, skulle den tala för att gå
tillbaka till det.

\needspace{10\baselineskip}
\section{Träffar som rör påståendet}

\Cref{tab:recap-flippat-kallor} listar de källor bilagan bygger på, med
hur de lästes, och \cref{tab:recap-flippat-motsager} de övriga träffar som
språkmodellen klassade som motsägande.  \Cref{tab:sok-recap-flippat}
listar alla träffar som rör påståendet: av de 845 är 20 citerade, 96
klassade som stödjande, 248 som nyanserande eller motsägande, 374 som
angränsande och 104 som felträffar, och 3 är dubbletter av citerade
källor; de angränsande, felträffarna och dubbletterna listas inte.

% Copied from research/evaluation/delutvärdering-2026/bilaga-flippat.tex
% (d536059) for the recap deck, bilaga F; \textcite in cells replaced by
% \citeauthor/\citeyear (no margin citations inside a table), label renamed.
\begingroup\footnotesize
\begin{longtable}{@{}>{\raggedright\arraybackslash}p{0.19\textwidth}%
    >{\raggedright\arraybackslash}p{0.16\textwidth}%
    >{\raggedright\arraybackslash}p{0.15\textwidth}%
    >{\raggedright\arraybackslash}p{0.38\textwidth}@{}}
  \caption{Källorna bilagan bygger på, efter roll. Läst anger om
    fulltexten eller bara sammanfattningen lästes.}
  \label{tab:recap-flippat-kallor}\\
  \toprule
  Källa & Slag & Läst & Bidrag \\
  \midrule
  \endfirsthead
  \multicolumn{4}{@{}l}{\emph{Tabell~\thetable{} (forts.)}}\\
  \toprule
  Källa & Slag & Läst & Bidrag \\
  \midrule
  \endhead
  \bottomrule
  \endlastfoot
  \multicolumn{4}{@{}l}{\emph{Metaanalyser av flippat klassrum mot
    föreläsning}}\\
  \citeauthor{LagSaele2019}~(\citeyear{LagSaele2019}) & metaanalys, 272 studier & fulltext &
    \(g = 0{,}35\); korrigerat 0,17--0,24; nöjdhet \(-0{,}03\) \\
  \citeauthor{vanAlten2019}~(\citeyear{vanAlten2019}) & metaanalys, 114 studier & fulltext &
    \(g = 0{,}36\); quiz \(+0{,}19\), kortad tid \(-0{,}26\); nöjdhet
    0,05 \\
  \citeauthor{Hew2021second}~(\citeyear{Hew2021second}) & metaanalys av 15 metaanalyser & fulltext &
    \(g = 0{,}45\), 0,37 efter korrigering; tecken på publiceringsbias \\
  \citeauthor{HewLo2018}~(\citeyear{HewLo2018}) & metaanalys, 28 studier, vård & fulltext &
    0,33; större effekt med quiz i början av passet \\
  \citeauthor{Cheng2019}~(\citeyear{Cheng2019}) & metaanalys, 55 publikationer & sammanfattning &
    \(g = 0{,}19\) \\
  \citeauthor{Strelan2020}~(\citeyear{Strelan2020}) & metaanalys, 198 studier & sammanfattning &
    \(g = 0{,}50\); IT 0,30 \\
  \citeauthor{Bredow2021}~(\citeyear{Bredow2021}) & metaanalys, 317 studier & sammanfattning &
    fördel i sju av åtta utfall, \(g\) 0,20--0,53 \\
  \citeauthor{Gillette2018}~(\citeyear{Gillette2018}) & metaanalys, 5 studier, farmaci &
    sammanfattning & ingen signifikant skillnad \\
  \midrule
  \multicolumn{4}{@{}l}{\emph{Metodgranskning och kontrollerade
    jämförelser}}\\
  \citeauthor{Kapur2022}~(\citeyear{Kapur2022}) & 46 metaanalyser och ny metaanalys &
    fulltext & effekter 0,19--2,29; fördelen krymper när den
    traditionella undervisningen är aktiv \\
  \citeauthor{Hew2021methodology}~(\citeyear{Hew2021methodology}) & granskning av 19 metaanalyser &
    sammanfattning & brister i sökning, biasbedömning och kontroll \\
  \citeauthor{Jensen2015}~(\citeyear{Jensen2015}) & kvasiexperiment, biologi & fulltext &
    aktivt flippat = aktivt icke-flippat \\
  \midrule
  \multicolumn{4}{@{}l}{\emph{Programmering, datavetenskap och
    matematik}}\\
  \citeauthor{AlmassriZaharudin2023}~(\citeyear{AlmassriZaharudin2023}) & metaanalys, 14 studier & fulltext &
    \(g = 0{,}56\), \(I^2 = 90\,\%\); svag metod \\
  \citeauthor{Tyler2018}~(\citeyear{Tyler2018}) & kursjämförelse över år & sammanfattning &
    ingen signifikant skillnad \\
  \citeauthor{Liu2026}~(\citeyear{Liu2026}) & jämförelse, statistisk programmering &
    sammanfattning & mindre lärandeökning med flippat \\
  \citeauthor{Giannakos2014}~(\citeyear{Giannakos2014}) & översikt, 32 artiklar & sammanfattning &
    efterlyser kontrollerade experiment \\
  \citeauthor{Jones2026}~(\citeyear{Jones2026}) & översikt, matematik & sammanfattning &
    stöd men inkonsekventa resultat \\
  \midrule
  \multicolumn{4}{@{}l}{\emph{Preferens och upplevt lärande}}\\
  \citeauthor{Deslauriers2019}~(\citeyear{Deslauriers2019}) & randomiserat försök, fysik & fulltext &
    lärde mer, upplevde mindre, föredrog föreläsning \\
  \citeauthor{Boedeker2025}~(\citeyear{Boedeker2025}) & randomiserat försök, medicin & fulltext &
    lärde mer och upplevde mer \\
  \citeauthor{Elliott2014}~(\citeyear{Elliott2014}) & enkätstudie, IT & sammanfattning & nöjda vid
    kursens slut \\
  \citeauthor{SoosaiRaj2018}~(\citeyear{SoosaiRaj2018}) & enkätstudie, inledande programmering &
    sammanfattning & föredrog livekodning framför kodning under passet \\
\end{longtable}
\endgroup


% Copied from research/evaluation/delutvärdering-2026/bilaga-flippat.tex
% (d536059) for the recap deck, bilaga F; label renamed.
\begingroup\footnotesize
\begin{longtable}{@{}>{\raggedright\arraybackslash}p{0.62\textwidth}%
    r>{\raggedright\arraybackslash}p{0.12\textwidth}r@{}}
  \caption{Övriga träffar som språkmodellen klassade som motsägande, med
    år, databas och modellens säkerhet. De lästes på
    sammanfattningsnivå men citeras inte; flera gäller andra upplägg än
    flippat klassrum, andra ämnen än programmering eller bara studenternas
    upplevelse.}
  \label{tab:recap-flippat-motsager}\\
  \toprule
  Titel & År & Databas & Säkerhet \\
  \midrule
  \endfirsthead
  \multicolumn{4}{@{}l}{\emph{Tabell~\thetable{} (forts.)}}\\
  \toprule
  Titel & År & Databas & Säkerhet \\
  \midrule
  \endhead
  \bottomrule
  \endlastfoot
  East versus west: The descendants of confucianism vs. evidence-based learning mainland Chinese learners in pursuit of western-based education in Singapore & 2013 & IEEE & 0,920 \\
  Inverted Classroom by Topic - A Study in Mathematics for Electrical Engineering Students & 2014 & OA & 0,980 \\
  Switching to blend-Ed: Effects of replacing the textbook with the browser in an introductory computer programming course & 2016 & IEEE & 0,960 \\
  A Flipped Classroom Model For Algorithm In College & 2017 & OA & 0,980 \\
  Flipped classroom approaches lead to no improvement in learning outcomes or student perceptions & 2017 & OA & 0,990 \\
  The flipped classroom: A learning model to increase student engagement not academic achievement & 2017 & OA & 0,990 \\
  ATTITUDES FOR IMPROVED LEARNING \& KNOWLEDGE IN ARCHETYPAL ENGINEERING COURSES & 2019 & WoS & 0,960 \\
  Flipping General Chemistry in Small Classes: Students’ Perception and Success & 2019 & OA & 0,970 \\
  Introducing an Inverted Classroom Concept for an Undergraduate Course in Power Electronics – Implementation, Evaluation and Experiences & 2019 & IEEE & 0,970 \\
  Undergraduates’ satisfaction and perceptions of learning outcomes across teacher- and learner-focused pedagogies & 2019 & Scopus & 0,970 \\
  Probing the Inverted Classroom: A Controlled Study of Teaching and Learning Outcomes in Undergraduate Engineering and Mathematics & 2020 & OA & 0,870 \\
  The Effect of Motivation on Learners' Performance and Satisfaction under Flipped Strategy in Discrete Math & 2020 & IEEE & 0,970 \\
  An Examination of Student Preferences and Learning Outcomes in Flipped Classroom with Online Videos & 2022 & OA & 0,990 \\
  Dependence of learning outcomes in flipped and lecture classrooms on review questions: A randomized controlled trial and observational study & 2022 & Scopus & 0,990 \\
  The Effectiveness of Using MOOC with E-FLIC Approach in Remote Learning as New Normal of Education to Support Basic Computer Science Learning & 2025 & IEEE & 0,980 \\
  Contribution of Active Learning Methods in Large-Group Teaching of Medical Microbiology & 2026 & Scopus & 0,960 \\
  Simulation-based education in an emergency medicine clerkship in Qatar: impact on academic performance and student perceptions & 2026 & Scopus & 0,970 \\
  The use of voiceovers and review games in teaching immunology & 2026 & Scopus & 0,970 \\
\end{longtable}
\endgroup


% Author's note: generated with `scholar sessions export
% flipped-classroom-learning -f table --bearing-only --lang sv --label
% flipped-traffar --theme ...`; copied from the report, label renamed.
% Copied from research/evaluation/delutvärdering-2026/litteratursokning/hits-table.tex
% (d536059) for the recap deck, bilaga F; only the label changed.
% --- scholar LaTeX fragment --------------------------------------------
% Generated by: scholar sessions export -f table
% Session: flipped-classroom-learning
% \input this file from the body of a document whose
% preamble loads:
%   \usepackage{longtable}
%   \usepackage{booktabs}
%   \usepackage{array}
% ----------------------------------------------------------------------

\begingroup\footnotesize
\begin{longtable}{@{}%
  >{\raggedright\arraybackslash}p{0.44\textwidth}c%
  >{\raggedright\arraybackslash}p{0.13\textwidth}%
  >{\raggedright\arraybackslash}p{0.26\textwidth}@{}}
\caption{Träffar som rör påståendet, flippat klassrum (20 citerade, 96 som stöder påståendet och 248 som kvalificerar eller motsäger det, av 845 unika träffar; de 374 angränsande och 104 felträffarna, 3 dubbletter eller andra versioner av citerade källor listas inte). Skälet till varje inkludering står i sista kolumnen; för maskinklassade rader anges modellens konfidens. Databaser: OpenAlex (OA; inte open access), IEEE Xplore (IEEE), Web of Science (WoS), Scopus.}\label{tab:sok-recap-flippat}\\
\toprule
Titel & År & Leverantör & Skäl \\
\midrule
\endfirsthead
\multicolumn{4}{@{}l}{\emph{\tablename~\thetable{} (forts.)}}\\
\toprule
Titel & År & Leverantör & Skäl \\
\midrule
\endhead
\midrule \multicolumn{4}{r@{}}{\emph{forts.\ på nästa sida}}\\
\endfoot
\bottomrule
\endlastfoot
A Comparison of Flipped Programming Classroom Models -- Results by Gender and Major & 2018 & IEEE & \emph{citerad}: programmering, datavetenskap eller matematik (refutes-claim) \\
A Meta-Analysis of Outcomes Comparing Flipped Classroom and Lecture & 2018 & OA & \emph{citerad}: metaanalys av flippat mot föreläsning (refutes-claim) \\
Active Versus Passive Learning in Large-Group Sessions in Medical School: A Randomized Cross-Over Trial Investigating Effects on Learning and the Feeling of Learning & 2025 & Scopus & \emph{citerad}: studentpreferens och upplevt lärande \\
Boundary conditions for the flipped classroom: Evidence from statistical programming & 2026 & WoS & \emph{citerad}: programmering, datavetenskap eller matematik (refutes-claim) \\
Do students like the flipped classroom? An investigation of student reaction to a flipped undergraduate IT course & 2014 & IEEE & \emph{citerad}: studentpreferens och upplevt lärande \\
Does the Flipped Classroom Improve Student Learning and Satisfaction? A Systematic Review and Meta-Analysis & 2019 & OA & \emph{citerad}: metaanalys av flippat mot föreläsning \\
Effectiveness of Flipped Classroom Pedagogy in Programming Education: A Meta-Analysis & 2023 & Scopus & \emph{citerad}: programmering, datavetenskap eller matematik \\
Effects of flipping the classroom on learning outcomes and satisfaction: A meta-analysis & 2019 & OA & \emph{citerad}: metaanalys av flippat mot föreläsning \\
Effects of the flipped classroom instructional strategy on students' learning outcomes: a meta-analysis & 2018 & OA & \emph{citerad}: metaanalys av flippat mot föreläsning \\
Fail, flip, fix, and feed -- Rethinking flipped learning: A review of meta-analyses and a subsequent meta-analysis & 2022 & OA & \emph{citerad}: metodkritik eller kontrollerad jämförelse \\
Flipped classroom improves student learning in health professions education: a meta-analysis & 2018 & WoS & \emph{citerad}: metaanalys av flippat mot föreläsning \\
Flipped learning in the mathematical sciences: insights, inconsistencies and gaps in the meta-analyses & 2026 & WoS & \emph{citerad}: programmering, datavetenskap eller matematik \\
Improvements from a Flipped Classroom May Simply Be the Fruits of Active Learning & 2015 & OA & \emph{citerad}: metodkritik eller kontrollerad jämförelse (refutes-claim) \\
Is More Active Always Better for Teaching Introductory Programming? & 2018 & IEEE & \emph{citerad}: studentpreferens och upplevt lärande \\
Measuring actual learning versus feeling of learning in response to being actively engaged in the classroom & 2019 & Scopus & \emph{citerad}: studentpreferens och upplevt lärande \\
Meta-analyses of flipped classroom studies: A review of methodology & 2021 & OA & \emph{citerad}: metodkritik eller kontrollerad jämförelse \\
On the use of flipped classroom across various disciplines: Insights from a second-order meta-analysis & 2021 & OA & \emph{citerad}: metaanalys av flippat mot föreläsning \\
Reviewing the flipped classroom research: Reflections for computer science education & 2014 & Scopus & \emph{citerad}: programmering, datavetenskap eller matematik \\
The flipped classroom: A meta-analysis of effects on student performance across disciplines and education levels & 2020 & OA & \emph{citerad}: metaanalys av flippat mot föreläsning \\
To Flip or Not to Flip? A Meta-Analysis of the Efficacy of Flipped Learning in Higher Education & 2021 & OA & \emph{citerad}: metaanalys av flippat mot föreläsning \\
A Flipped Classroom Approach to Teaching Concurrency and Parallelism & 2016 & IEEE & stöder påståendet (0.97) \\
A novel extended flipped classroom model helps dental undergraduates grow into dentists & 2026 & Scopus & stöder påståendet (0.95) \\
A study of inverted classroom pedagogy in computer science teaching & 2015 & OA & stöder påståendet (0.98) \\
A systematic review and meta-analysis of flipped learning among university students in Korea: Self-directed learning, learning motivation, efficacy, and learning achievement & 2021 & OA & stöder påståendet (1.00) \\
Academic Achievement, Student Experience, and Behavioral Engagement: A Mixed-Methods and Learning Analytics Investigation of the Blended-Flipped Rotation Model & 2026 & Scopus & stöder påståendet (0.98) \\
Active Learning Strategies in Computer Science Education: A Systematic Review & 2024 & Scopus & stöder påståendet (0.88) \\
Active learning through flipped classroom in mechanical engineering: improving students' perception of learning and performance & 2021 & OA & stöder påståendet (0.99) \\
An analysis of flip-classroom pedagogy in first year undergraduate mathematics for computing & 2014 & IEEE & stöder påståendet (0.98) \\
An empirical analysis of the flipped classroom about information literacy education in Chinese university based on SPSS 19.0 & 2021 & IEEE & stöder påståendet (0.90) \\
Analysis of Postgraduate Theses on Flipped Learning: A Meta-Analysis Study & 2023 & OA & stöder påståendet (0.88) \\
Application of the BOPPPS model combined with flipped classroom in clinical practice competency training for gastroenterology residents in standardized residency programs & 2026 & Scopus & stöder påståendet (0.96) \\
Assessing the Impact of the LearnOOP Tool on Enhancing the Understanding of Object-Oriented Programming-An Experimental Study Across Three Semesters & 2026 & IEEE & stöder påståendet (0.86) \\
AURA-Based Ranking of Teaching Strategies for Business Law: Toward Data-Informed Pedagogy & 2025 & IEEE & stöder påståendet (0.75) \\
Blended peer-led research curriculum with AI integration improves postgraduate students' academic performance and satisfaction: a quasi-experimental mixed-methods study & 2026 & Scopus & stöder påståendet (0.96) \\
Blending Inverted Lectures and Laboratory Experiments to Improve Learning in an Introductory Course in Digital Systems & 2020 & IEEE & stöder påståendet (0.98) \\
Boosting Engineering Students' Trigonometry and Solid Mensuration Skills Through Flipped Learning & 2026 & Scopus & stöder påståendet (0.99) \\
Changing the landscape of medical anatomy education: impact of flipped classroom pedagogy on learning outcomes and perceptions of medical students in practical gastrointestinal anatomy & 2026 & Scopus & stöder påståendet (0.99) \\
College students' cognitive learning outcomes in flipped classroom instruction: a meta-analysis of the empirical literature & 2019 & OA & stöder påståendet (0.97) \\
Comparative Evaluation of Traditional and Jigsaw-based Learning Approach in the Understanding of Class I Cavity Design among Dental Undergraduates & 2026 & Scopus & stöder påståendet (0.72) \\
Comparing student achievement in traditional learning with a combination of blended and flipped learning & 2020 & OA & stöder påståendet (0.99) \\
Comparison of the BOPPPS model combined with rain classroom teaching and traditional method on nephrology education for fifth-year undergraduates & 2026 & Scopus & stöder påståendet (0.91) \\
Comprehensive structured review of implementing flipped classroom approaches in education & 2025 & Scopus & stöder påståendet (0.88) \\
Constructing and Validating a Flipped Classroom Framework for College Physical Education & 2026 & Scopus & stöder påståendet (0.97) \\
Design and Assessment of a Hybrid Teaching Model for Teaching Internet of Things Hardware Technology & 2026 & Scopus, WoS & stöder påståendet (0.86) \\
Does the Flipped Classroom Lead to Increased Gains of Learning Outcomes in ESL/EFL Contexts? & 2016 & OA & stöder påståendet (0.98) \\
Effectiveness of Camtasia-supported flipped learning on improving self-learning achievement among female students at the college of basic education & 2026 & WoS & stöder påståendet (0.94) \\
Effectiveness of flipped classroom in pharmacy education - a meta-analysis & 2023 & WoS & stöder påståendet (0.99) \\
Effectiveness of flipped classroom on learning and satisfaction of nursing students in physiology laboratory education: a quasi-experimental study & 2026 & Scopus & stöder påståendet (0.97) \\
Effectiveness of flipped classrooms in Chinese baccalaureate nursing education: A meta-analysis of randomized controlled trials & 2017 & OA & stöder påståendet (0.97) \\
Effectiveness of flipped classrooms in nursing education: Systematic review and meta-analysis & 2017 & OA & stöder påståendet (0.93) \\
Effects of using flipped classroom learning in object-oriented analysis and design course & 2017 & IEEE & stöder påståendet (0.86) \\
Enhancing Anatomy Education: Snowball Techniques vs. Lecture-Based Learning & 2026 & WoS & stöder påståendet (0.82) \\
Enhancing flipped classroom model implementation & 2017 & IEEE & stöder påståendet (0.83) \\
Enhancing programming education through AI-driven tools and innovative pedagogical models: a systematic study on student engagement, retention, and self-regulation & 2026 & Scopus & stöder påståendet (0.76) \\
Enhancing Statistical Thinking in Higher Education Through Pedagogically Designed Use of Interactive Whiteboards & 2026 & WoS & stöder påståendet (0.96) \\
Evaluating an ICT-Mediated active learning framework: A quasi-experimental study in ukrainian higher education & 2025 & WoS & stöder påståendet (0.99) \\
Evaluating the effectiveness of flipped learning in an upper-division undergraduate electrical engineering course & 2024 & OA & stöder påståendet (0.96) \\
Exploring the Effects of Online Academic Help-Seeking and Flipped Learning on Improving Students' Learning. & 2017 & OA & stöder påståendet (0.94) \\
Flipped Classroom Applied to Software Architecture Teaching & 2020 & IEEE & stöder påståendet (0.91) \\
Flipped classroom combined with case-based learning enhances ventilator management training for non-anesthesiology medical undergraduates: a randomized controlled study & 2026 & Scopus & stöder påståendet (0.96) \\
Flipped Classroom for Linear Algebra at Undergraduate Level & 2024 & IEEE & stöder påståendet (0.91) \\
Flipped classroom improves nursing students' theoretical learning in China: A meta-analysis & 2020 & OA & stöder påståendet (0.99) \\
Flipped classroom improves student learning outcome in Chinese pharmacy education: A systematic review and meta-analysis & 2022 & OA & stöder påståendet (0.99) \\
Flipped classroom in accounting courses: A systematic review & 2022 & OA & stöder påståendet (0.94) \\
Flipped classroom integrated with team-based learning enhances surgical skills in ophthalmology residency training: a randomized controlled trial & 2026 & WoS & stöder påståendet (0.98) \\
Flipped Classroom Method Improves Anatomy and Histology Lab Education: Evaluating the Effects on Satisfaction and Learning in Dental Students & 2026 & WoS & stöder påståendet (0.99) \\
Flipped Classrooms in Higher Education: A Meta-Analysis of Their Impact on Mathematics Performance & 2025 & WoS & stöder påståendet (0.98) \\
Flipped Learning Outcomes among Undergraduate Medical Students, Systematic Review and Met analysis, 2022 & 2022 & OA & stöder påståendet (0.78) \\
Flipped pedagogy: Strategies and technologies in chemistry education & 2021 & Scopus & stöder påståendet (0.72) \\
Flipped-classroom Based Teaching Method for Data Structures and Algorithms & 2025 & WoS & stöder påståendet (0.96) \\
Flipping the Culture, not just the Classroom: Leadership Lessons from a Residency Curriculum Redesign & 2026 & WoS & stöder påståendet (0.91) \\
Flipping Web Programming Class: Student's Perception and Performance & 2019 & IEEE & stöder påståendet (0.98) \\
Impact of blended learning on learning outcomes in the public healthcare education course: a review of flipped classroom with team-based learning & 2021 & OA & stöder påståendet (0.98) \\
Impact of Flipped Classroom Instruction on Brain-Mediated Motor Skill Performance in University Students: A Systematic Review and Meta-Analysis & 2025 & WoS & stöder påståendet (0.94) \\
Impact of the Pedagogy of Successive Questioning on Chemistry Learning & 2025 & WoS & stöder påståendet (0.97) \\
Implementing the Flipped Classroom in Teacher Education: Evidence from Turkey. & 2017 & OA & stöder påståendet (0.84) \\
Improving student learning in an introductory programming course using flipped classroom and competency framework & 2017 & IEEE & stöder påståendet (0.99) \\
Increasing Students' Responsibility and Learning Outcomes Using Partial Flipped Classroom in a Language Processors Course & 2020 & IEEE & stöder påståendet (0.99) \\
Innovative Pedagogy and Design-Based Research on Flipped Learning in Higher Education & 2021 & OA & stöder påståendet (0.90) \\
Integration of Flipped Classroom and Problem Based Learning Model and its Implementation in University Programming Course & 2019 & IEEE & stöder påståendet (0.99) \\
Is the flipped classroom more effective than the traditional classroom in clinical medical education: a systematic review and meta-analysis & 2024 & Scopus & stöder påståendet (0.93) \\
Learning with technology in physiotherapy education: Design, implementation and evaluation of a flipped classroom teaching approach & 2019 & Scopus & stöder påståendet (0.94) \\
Meta-analysis of flipped classroom model to promote mathematical higher-order thinking & 2026 & Scopus & stöder påståendet (0.88) \\
Modern Didactic Formats in Surgery: A Systematic Review & 2023 & WoS & stöder påståendet (0.96) \\
Motivational Teaching Techniques in Secondary and Higher Education: A Systematic Review of Active Learning Methodologies & 2025 & Scopus & stöder påståendet (0.68) \\
Peer Role-Play Integration in Flipped Classroom Education Enhances Medical Students' Clinical Reasoning in Child Abuse Management: A Quasi-Experimental Study & 2026 & WoS & stöder påståendet (0.93) \\
Perception of case-based collaborative learning in pathology through panel discussion using concept mapping: a quasi-experimental crossover study in phase II MBBS students & 2026 & Scopus & stöder påståendet (0.91) \\
Performance and Perception in the Flipped Learning Model: An Initial Approach to Evaluate the Effectiveness of a New Teaching Methodology in a General Science Classroom & 2016 & OA & stöder påståendet (0.93) \\
Pharmacy Student Engagement, Performance, and Perception in a Flipped Satellite Classroom & 2013 & OA & stöder påståendet (0.91) \\
Probing the Flipped Learning Literature in Social Sciences and Humanities Education & 2022 & OA & stöder påståendet (0.95) \\
Research to Practice in Computer Programming Course using Flipped Classroom & 2022 & IEEE & stöder påståendet (0.86) \\
Student learning and perceptions in a flipped linear algebra course & 2013 & OA & stöder påståendet (0.96) \\
Student Learning Performance and Satisfaction With a Flipped Classroom in Undergraduate Dental Pharmacology Education & 2026 & Scopus & stöder påståendet (0.99) \\
The Effect of Flipped Classroom for Writing: A Meta Analysis & 2026 & WoS, Scopus & stöder påståendet (0.86) \\
The Effect of Flipped Learning Approach on Academic Achievement: A Meta-Analysis Study & 2018 & OA & stöder påståendet (0.99) \\
The Effect of Flipped Learning on Students' Academic Achievement: A Meta-Analysis Study & 2019 & OA & stöder påståendet (0.98) \\
The effectiveness of the flipped learning model in higher education: A meta-analysis study & 2024 & OA & stöder påståendet (0.76) \\
The effects of flipped classroom versus lecture-based learning on academic performance and perceptions of classroom environment among nursing students: a randomized controlled trial & 2026 & Scopus & stöder påståendet (0.91) \\
The effects of flipped classrooms on undergraduate pharmaceutical marketing learning: A clustered randomized controlled study & 2019 & OA & stöder påståendet (1.00) \\
The flipped classroom in engineering education: A survey of the research & 2015 & IEEE & stöder påståendet (0.98) \\
The flipped classroom model and its influence on academic performance: A systematic review & 2026 & WoS & stöder påståendet (0.93) \\
The flipped classroom stimulates greater learning and is a modern 21st century approach to teaching today's undergraduates & 2015 & OA & stöder påståendet (0.91) \\
The Impact of Flipped Learning on Achievement and Attitudes In Higher Education & 2014 & OA & stöder påståendet (0.97) \\
The influences of Flipped Classroom: A meta analysis & 2014 & IEEE & stöder påståendet (0.79) \\
THE META-ANALYSIS OF THE STUDIES ABOUT THE EFFECTS OF FLIPPED LEARNING ON STUDENTS' ACHIEVEMENT & 2020 & OA & stöder påståendet (0.97) \\
The relative effect of team-based learning on motivation and learning: A self-determination theory perspective & 2017 & Scopus & stöder påståendet (0.91) \\
The Study of Programming Language Learning by Applying Flipped Classroom & 2018 & IEEE & stöder påståendet (0.91) \\
The trends and outcomes of flipped learning research between 2012 and 2018: A descriptive content analysis & 2021 & OA & stöder påståendet (0.92) \\
Transdisciplinarity and Team-Based Learning: Strategies for an Introductory Programming Course & 2024 & IEEE & stöder påståendet (0.86) \\
Transforming Education: Enhancing Student Performance and Satisfaction through the Flipped Classroom Method & 2024 & OA & stöder påståendet (0.93) \\
Using a Partially Flipped Classroom and Gamification to Improve Student Performance in a First-Year Electronic Circuits Course & 2024 & IEEE & stöder påståendet (0.99) \\
Using flipped classroom, peer discussion, and just-in-time teaching to increase learning in a programming course & 2015 & IEEE & stöder påståendet (0.98) \\
Visual Learning in Electrocardiography Training for Medical Residents: Comparative Intervention Study & 2025 & WoS & stöder påståendet (0.82) \\
Vodcasts and Active-Learning Exercises in a ``Flipped Classroom'' Model of a Renal Pharmacotherapy Module & 2012 & OA & stöder påståendet (0.99) \\
WIP: Novel Pre-Lecture Technique for Accelerated Student Engagement in Circuit Analysis Courses & 2025 & IEEE & stöder påståendet (0.91) \\
Would Flipped Classroom be My Approach in Teaching Computing Courses: Literature Review & 2021 & Scopus & stöder påståendet (0.91) \\
 & 2023 & OA, Scopus & kvalificerar eller motsäger påståendet (0.94) \\
``Flipping'' The Classroom To Explore Active Learning In A Large Undergraduate Course & 2020 & OA & kvalificerar eller motsäger påståendet (0.79) \\
A comparison of student evaluations of teaching and learning in the inverted classroom model versus traditional lectures in dental education & 2026 & OA, Scopus & kvalificerar eller motsäger påståendet (0.91) \\
A Flipped Classroom Model For Algorithm In College & 2017 & OA & kvalificerar eller motsäger påståendet (0.98) (refutes-claim) \\
A flipped classroom model in higher education: a review of the evidence across disciplines & 2019 & OA & kvalificerar eller motsäger påståendet (0.94) \\
A Meta Analysis on Effects of Flipped Learning in Korea & 2018 & OA & kvalificerar eller motsäger påståendet (0.78) \\
A Meta-Analysis of Learning and Moderator Variables in Flipped Education & 2022 & OA & kvalificerar eller motsäger påståendet (0.99) \\
A meta-analysis on flipped learning: Conditions for successful application and future research direction & 2016 & OA & kvalificerar eller motsäger påståendet (0.95) \\
A MULTIDISCIPLINARY UMBRELLA REVIEW ON FLIPPED LEARNING OUTCOMES & 2025 & OA & kvalificerar eller motsäger påståendet (0.98) \\
A novel integration of online and flipped classroom instructional models in public health higher education & 2014 & OA & kvalificerar eller motsäger påståendet (0.96) \\
A review of flipped learning in innovative math education & 2023 & OA & kvalificerar eller motsäger påståendet (0.79) \\
A structured blended learning approach for medical imaging technology interns: a prospective randomized comparative study & 2026 & Scopus & kvalificerar eller motsäger påståendet (0.91) \\
A Study of Students Perspectives on Different Pedagogical Practices for Remote Digital Signal Processing Courses & 2021 & IEEE & kvalificerar eller motsäger påståendet (0.89) \\
A systematic review of flipped classroom approaches in language learning & 2024 & Scopus & kvalificerar eller motsäger påståendet (0.80) \\
A systematic review of flipped classroom empirical evidence from different fields: what are the gaps and future trends? & 2019 & OA & kvalificerar eller motsäger påståendet (0.94) \\
A Systematic Review of Self-Regulated Learning in Flipped Classrooms: Key Findings, Measurement Methods, and Potential Directions & 2022 & IEEE & kvalificerar eller motsäger påståendet (0.94) \\
A systematic review of the use of gamification in flipped learning & 2021 & OA & kvalificerar eller motsäger påståendet (0.84) \\
A Systematic Review of Theoretical Foundations and Learning Effects in Gamified Flipped Classroom Research & 2025 & WoS & kvalificerar eller motsäger påståendet (0.90) \\
Academic outcomes of flipped classroom learning: a metaanalysis & 2018 & OA & kvalificerar eller motsäger påståendet (0.99) \\
Academic Performance and Perceptions of Undergraduate Medical Students in Case-Based Learning Compared to Other Teaching Strategies: A Systematic Review with Meta-Analysis & 2023 & WoS & kvalificerar eller motsäger påståendet (0.82) \\
Active learning in diabetes education for health professions students: a narrative synthesis of effectiveness and implementation considerations & 2026 & WoS & kvalificerar eller motsäger påståendet (0.92) \\
Active Learning in Flipped Life Science Courses Promotes Development of Critical Thinking Skills & 2018 & OA & kvalificerar eller motsäger påståendet (0.82) \\
Active Learning Strategies in Calculus Education: A Comparative Study Between French and Brazilian Engineering Students & 2026 & WoS & kvalificerar eller motsäger påståendet (0.82) \\
Active learning under time constraints: Working students, institutional conditions and implementation fit in a public business school in Brazil & 2026 & WoS & kvalificerar eller motsäger påståendet (0.82) \\
Active learning, cooperative active learning, and passive learning methods in an accounting information systems course & 2017 & Scopus & kvalificerar eller motsäger påståendet (0.99) \\
Active methodologies in medical education: between pedagogical innovation and the persistence of the lecture-based class. & 2026 & WoS & kvalificerar eller motsäger påståendet (0.83) \\
Adapting Dental Education for the Gen Z: An Overview of Active Learning Strategies & 2026 & WoS & kvalificerar eller motsäger påståendet (0.83) \\
Addressing the challenges and impact of flipped classroom model on the performance and perception of health professional students: an umbrella review & 2026 & Scopus & kvalificerar eller motsäger påståendet (0.96) \\
Advances in medical education and practice: student perceptions of the flipped classroom & 2017 & OA & kvalificerar eller motsäger påståendet (0.94) \\
An Examination of Student Preferences and Learning Outcomes in Flipped Classroom with Online Videos & 2022 & OA & kvalificerar eller motsäger påståendet (0.99) (refutes-claim) \\
Analyses of possibilities of Flipped Classroom in Teaching Computer Science Courses & 2019 & IEEE & kvalificerar eller motsäger påståendet (0.83) \\
Analysis of the Studies Concerning Flipped Learning Model: A Comparative Meta-Synthesis Study & 2018 & OA & kvalificerar eller motsäger påståendet (0.83) \\
Application of Emerging Teaching Models in Dental Education: A Systematic Review and Meta-Analysis & 2024 & Scopus & kvalificerar eller motsäger påståendet (0.89) \\
Application of flipped classroom in surgical education: a systematic review and meta-analysis & 2026 & WoS & kvalificerar eller motsäger påståendet (0.94) \\
Applying Flipped Classroom with a Remote Laboratory for Photovoltaic Energy Efficiency Experimentation & 2026 & IEEE & kvalificerar eller motsäger påståendet (0.78) \\
Applying Peer Instruction to Computer Science Students Using Non-native Language: A Study with Undergraduate Students & 2021 & IEEE & kvalificerar eller motsäger påståendet (0.81) \\
Are first year students ready for a flipped classroom? A case for a flipped learning continuum & 2019 & OA & kvalificerar eller motsäger påståendet (0.97) \\
Are online lectures effective? Evidence from a flipped classroom pilot in UK higher education & 2026 & OA & kvalificerar eller motsäger påståendet (0.94) \\
Assessing Learning Outcomes and Evaluating Graduate Student Perceptions of a Flipped Classroom & 2016 & OA & kvalificerar eller motsäger påståendet (0.95) \\
Between lecture and flipped: The use of ``labs'' in the public administration classroom & 2019 & Scopus & kvalificerar eller motsäger påståendet (0.88) \\
Case-Based Instruction in STEM: Analysis of Student Confidence & 2016 & IEEE & kvalificerar eller motsäger påståendet (0.82) \\
Case-Based VR-FL Versus Traditional Case Teaching: An Evaluation of Learning Achievement and Self-Efficacy among Nursing Students & 2025 & IEEE & kvalificerar eller motsäger påståendet (0.91) \\
Characterizing Engineering Learners' Preferences for Active and Passive Learning Methods & 2018 & IEEE & kvalificerar eller motsäger påståendet (0.94) \\
Comparing Flipped Classroom and Traditional Lecture in Medical Histology Among Chinese Medical Students: Learning Outcomes and Student Feedback & 2026 & WoS & kvalificerar eller motsäger påståendet (0.80) \\
Comparing the Effectiveness of an Inverted Classroom to a Traditional Classroom in an Upper-Division Engineering Course & 2013 & OA & kvalificerar eller motsäger påståendet (0.99) \\
Comparison of Pharmaceutical Calculations Learning Outcomes Achieved Within a Traditional Lecture or Flipped Classroom Andragogy & 2017 & OA & kvalificerar eller motsäger påståendet (0.99) \\
Comparison of the effects of different teaching methods on the effectiveness of teaching neurology in China: a bayesian network meta-analysis and systematic review & 2024 & WoS & kvalificerar eller motsäger påståendet (0.96) \\
Contribution of Active Learning Methods in Large-Group Teaching of Medical Microbiology & 2026 & Scopus & kvalificerar eller motsäger påståendet (0.96) (refutes-claim) \\
Data Structures in Flipped Classroom: Students' Effort and Preference & 2015 & IEEE & kvalificerar eller motsäger påståendet (0.95) \\
Deaf and Hard of Hearing Students' Perceptions of the Flipped Classroom Strategy in an Undergraduate Education Course & 2019 & OA & kvalificerar eller motsäger påståendet (0.88) \\
Dependence of learning outcomes in flipped and lecture classrooms on review questions: A randomized controlled trial and observational study & 2022 & Scopus & kvalificerar eller motsäger påståendet (0.99) (refutes-claim) \\
Design, implementation, and evaluation of an inverted (flipped) classroom model economics for sustainable education course & 2018 & OA & kvalificerar eller motsäger påståendet (0.88) \\
Discrepancies between Student Perception and Achievement of Learning Outcomes in a Flipped Classroom & 2016 & OA & kvalificerar eller motsäger påståendet (0.96) \\
Does a flipped classroom model work in mathematics education? A meta-analysis & 2023 & Scopus & kvalificerar eller motsäger påståendet (0.98) \\
Does flipped classroom improve students' evaluations of teaching quality? & 2026 & WoS & kvalificerar eller motsäger påståendet (0.98) \\
East versus west: The descendants of confucianism vs. evidence-based learning mainland Chinese learners in pursuit of western-based education in Singapore & 2013 & IEEE & kvalificerar eller motsäger påståendet (0.92) (refutes-claim) \\
Educational methods for inverted-lecture computer science classrooms to overcome common barriers to STEM student success & 2016 & IEEE & kvalificerar eller motsäger påståendet (0.94) \\
Effectiveness Evaluation of Integrating Chinese Classical Music Narrative Into the Flipped Classroom & 2026 & Scopus & kvalificerar eller motsäger påståendet (0.91) \\
Effectiveness of a case-based learning combined with flipped classroom approach in undergraduate pulmonary imaging education: a prospective randomized controlled study & 2026 & Scopus & kvalificerar eller motsäger påståendet (0.92) \\
Effectiveness of artificial intelligence chatbot-assisted flipped learning on patient handover for nursing students: A mixed-methods study & 2026 & WoS & kvalificerar eller motsäger påståendet (0.86) \\
Effectiveness of Different Educational Approaches in Improving Comprehensive Clinical Competence Among Medical Students: A Systematic Review and Meta-Analysis & 2026 & Scopus & kvalificerar eller motsäger påståendet (0.86) \\
Effectiveness of Flipped Classroom and Traditional Didactic Classroom in Paediatric Dentistry Education---A Systematic Review With Meta-Analysis of Randomised Controlled Trials & 2026 & Scopus & kvalificerar eller motsäger påståendet (0.97) \\
Effectiveness of flipped classroom practices in teaching of science: a mixed research synthesis & 2023 & Scopus & kvalificerar eller motsäger påståendet (0.94) \\
Effectiveness of flipped classrooms in Chinese students of clinical medicine major undergoing clinical practice: a meta-analysis & 2025 & Scopus & kvalificerar eller motsäger påståendet (0.97) \\
Effectiveness of Flipped Learning in Nursing Education: A Meta-Analysis Based on 12 Studies & 2025 & OA & kvalificerar eller motsäger påståendet (0.99) \\
Effectiveness of Geogebra Integration into Flipped Classroom (GFC) on Students Mathematics Skills: A Meta-Analysis Study & 2024 & OA & kvalificerar eller motsäger påståendet (0.90) \\
Effectiveness of Innovative Teaching Methods Versus Traditional Lectures in Dental Education: A Systematic Review and Meta-Analysis & 2026 & Scopus & kvalificerar eller motsäger påståendet (0.90) \\
Effectiveness of the flipped classroom in medical education: a systematic review and meta-analysis & 2026 & Scopus & kvalificerar eller motsäger påståendet (0.98) \\
Effectiveness of the Flipped Classroom in post-pandemic universities: cognitive impact and competence development from a systematic review PRISMA (2021-2024) & 2026 & WoS & kvalificerar eller motsäger påståendet (0.98) \\
Effectiveness of the Flipped Classroom Strategy in Learning Outcomes (Bibliometric Study) & 2017 & OA & kvalificerar eller motsäger påståendet (0.84) \\
Effects of an Integrated ARCS-Flipped Model on Technology-Enhanced Chinese Listening and Motivation & 2026 & WoS & kvalificerar eller motsäger påståendet (0.98) \\
Effects of flipped language classrooms on learning outcomes in higher education: A Bayesian meta-analysis & 2023 & WoS & kvalificerar eller motsäger påståendet (0.99) \\
Effects of Flipped Learning on Language Learning Outcomes: A Meta-Analysis investigating Moderators & 2025 & Scopus & kvalificerar eller motsäger påståendet (0.98) \\
Effects of Flipped Learning on Learning Outcomes and Satisfaction: A Meta-Analysis -Supplementary material & 2020 & OA & kvalificerar eller motsäger påståendet (0.98) \\
Effects of problem-based learning and flipped classroom educational methods on patient safety knowledge, attitudes and skills among nursing students: A systematic review and network meta-analysis & 2026 & Scopus & kvalificerar eller motsäger påståendet (0.96) \\
Effects of the BOPPPS model combined with case-based learning on knowledge acquisition and learning engagement in undergraduate nursing students: a multi-cohort quasi-experimental study & 2026 & Scopus & kvalificerar eller motsäger påståendet (0.87) \\
Efficacy of a Flipped Method in an Undergraduate Class at IIT Bombay & 2016 & IEEE & kvalificerar eller motsäger påståendet (0.96) \\
Efficacy of replacing the lecture with a SKILL in engineering science courses & 2015 & Scopus & kvalificerar eller motsäger påståendet (0.68) \\
EFL Students' Perceptions on Flipped Learning in an English Grammar Class & 2023 & OA & kvalificerar eller motsäger påståendet (0.87) \\
Engaging students in education for sustainable development: The benefits of active learning, reflective practices and flipped classroom pedagogies & 2021 & OA & kvalificerar eller motsäger påståendet (0.89) \\
Engineering Students' Perceptions of a Flipped English-Medium Instruction Course: Insights from a Taiwanese University & 2026 & Scopus & kvalificerar eller motsäger påståendet (0.68) \\
Engineering students' preferences for active and non-active learning strategies in the classroom & 2021 & IEEE & kvalificerar eller motsäger påståendet (0.93) \\
Enhancing CASEL 5 Social and Emotional Competencies in Engineering Students through Flipped and Gamified Blended Learning: A Longitudinal Study & 2026 & Scopus & kvalificerar eller motsäger påståendet (0.91) \\
Enhancing competency and self-directed learning in anesthesiology residency: an outcome-based education model integrating online-offline hybrid teaching and mind mapping: a randomized controlled trial & 2026 & WoS & kvalificerar eller motsäger påståendet (0.94) \\
Enhancing Conceptual Learning in Chemical Engineering and Nanotechnology Through the Flipped-Classroom Model and AI-Curated Modules & 2026 & Scopus & kvalificerar eller motsäger påståendet (0.96) \\
Enhancing human physiology education through active learning and digital tools: a modern pedagogical approach & 2025 & OA & kvalificerar eller motsäger påståendet (0.92) \\
Enhancing learning performance and soft skills development through body painting in human anatomy education & 2026 & Scopus & kvalificerar eller motsäger påståendet (0.91) \\
Enhancing physics education through flipped problem-based learning (PBL): improving comprehension and problem-solving skills in undergraduate students & 2026 & Scopus & kvalificerar eller motsäger påståendet (0.99) \\
Enhancing the Fourier Transform Learning Capabilities of Second-Year Electrical Engineering Students through Flipped Classroom Approach & 2025 & IEEE & kvalificerar eller motsäger påståendet (0.90) \\
Evaluating the Effectiveness of Flipped Lesson Models in Physical Education: A Systematic Review of Student Engagement and Learning Outcomes & 2026 & OA & kvalificerar eller motsäger påståendet (0.87) \\
Evaluating the effectiveness of technology-assisted learning versus traditional methods in histology education: A systematic review and meta-analysis & 2026 & Scopus & kvalificerar eller motsäger påståendet (0.92) \\
Evaluating The Flipped Classroom Approach in Computer Science Curricula & 2025 & Scopus & kvalificerar eller motsäger påståendet (0.85) \\
Evaluation of a new engineering teaching approach through a quasi-experimental study and student perceptions & 2026 & Scopus & kvalificerar eller motsäger påståendet (0.94) \\
Evaluation the effectiveness of the multimodal situational blended teaching model in higher vocational pediatric nursing education: a quasi-experimental study & 2026 & Scopus & kvalificerar eller motsäger påståendet (0.78) \\
Examining Effects of Science Flipped Classrooms A Meta-analysis & 2022 & IEEE & kvalificerar eller motsäger påståendet (0.99) \\
Exploring flipped classrooms in undergraduate anatomy and physiology courses: a systematic review & 2026 & Scopus & kvalificerar eller motsäger påståendet (0.98) \\
Exploring flipped classrooms in undergraduate nursing and health science: A systematic review & 2022 & WoS & kvalificerar eller motsäger påståendet (0.98) \\
Exploring Flipped Learning in an Introductory Programming Module: A Literature Review & 2024 & Scopus & kvalificerar eller motsäger påståendet (0.94) \\
Exploring how medical students learn during clinical rotations: a pilot study with a mobile application & 2019 & Scopus & kvalificerar eller motsäger påståendet (0.84) \\
Exploring the Effectiveness of Flipped Classroom on STEM Student Achievement: A Meta-analysis & 2023 & OA & kvalificerar eller motsäger påståendet (1.00) \\
Exploring the effectiveness of flipped learning in university education & 2023 & OA & kvalificerar eller motsäger påståendet (0.82) \\
Exploring the pedagogical design features of the flipped classroom in undergraduate nursing education: a systematic review & 2021 & OA & kvalificerar eller motsäger påståendet (0.95) \\
Flexibility vs. disconnection: A qualitative inquiry into student experiences in a flipped learning environment & 2026 & Scopus & kvalificerar eller motsäger påståendet (0.93) \\
Flipped but engaged? A systematic review of student engagement in higher education flipped classrooms incorporating the agentic dimension & 2026 & Scopus & kvalificerar eller motsäger påståendet (0.87) \\
Flipped classroom -- a qualitative study of medical students & 2026 & OA & kvalificerar eller motsäger påståendet (0.91) \\
Flipped classroom and learning outcomes: a second-order meta-analysis & 2026 & Scopus & kvalificerar eller motsäger påståendet (0.99) \\
Flipped classroom approach in financial accounting & 2025 & Scopus & kvalificerar eller motsäger påståendet (0.95) \\
Flipped classroom approach of language education: a systematic review & 2025 & Scopus & kvalificerar eller motsäger påståendet (0.84) \\
Flipped Classroom Approaches in Computer Programming Courses in Japan & 2020 & IEEE & kvalificerar eller motsäger påståendet (0.91) \\
Flipped Classroom Design as a Driver of Digital Transformation and Sustainable Education in Higher Education: A Systematic Review of Reviews & 2026 & Scopus & kvalificerar eller motsäger påståendet (0.99) \\
Flipped classroom goes sideways: reflections on active learning methodologies & 2022 & OA & kvalificerar eller motsäger påståendet (0.88) \\
Flipped classroom improves academic achievement, learning retention and attitude towards course: a meta-analysis & 2021 & OA & kvalificerar eller motsäger påståendet (0.99) \\
Flipped classroom in higher education: a systematic literature review and research challenges & 2023 & OA & kvalificerar eller motsäger påståendet (0.89) \\
Flipped classroom in technology courses - impact on personal efficacy and perception based on learning style preferences & 2017 & IEEE & kvalificerar eller motsäger påståendet (0.87) \\
Flipped Classroom Increases Achievement of Student Learning Outcomes & 2018 & OA & kvalificerar eller motsäger påståendet (0.89) \\
Flipped Classroom Strategy to Help Underachievers in Java Programming & 2018 & IEEE & kvalificerar eller motsäger påståendet (0.93) \\
Flipped classroom technology as a multicultural inclusion strategy for international medical students in Vietnamese political theory education & 2026 & WoS & kvalificerar eller motsäger påståendet (0.88) \\
Flipped Classroom, Flipped Teaching and Flipped Learning in the Foreign/Second Language Post--Secondary Classroom & 2018 & OA & kvalificerar eller motsäger påståendet (0.72) \\
Flipped classroom: a challenge to university education & 2020 & WoS & kvalificerar eller motsäger påståendet (0.95) \\
Flipped classrooms and misconceptions in biology: a systematic review & 2026 & Scopus & kvalificerar eller motsäger påståendet (0.70) \\
Flipped classrooms and student learning: not just surface gains & 2016 & OA & kvalificerar eller motsäger påståendet (0.89) \\
Flipped classrooms in pharmacy education: A systematic review & 2023 & OA & kvalificerar eller motsäger påståendet (0.93) \\
Flipped interpreting classroom: flipping approaches, student perceptions, and design considerations & 2017 & OA & kvalificerar eller motsäger påståendet (0.85) \\
Flipped learning : A literature review & 2022 & OA & kvalificerar eller motsäger påståendet (0.88) \\
Flipped Learning in Fashion Design with AI: Evaluating Student Attitudes, Learning Reflections, and AI-Mediated Lecturer Presence & 2026 & IEEE & kvalificerar eller motsäger påståendet (0.78) \\
Flipped Learning in Flipped Classrooms: A New Pathway to Prepare Future Special Educators & 2019 & OA & kvalificerar eller motsäger påståendet (0.85) \\
Flipped Learning in Foreign Language Learning in Higher Education: Analysis of Effectiveness and Moderator Variables & 2025 & OA & kvalificerar eller motsäger påståendet (0.99) \\
Flipped learning in higher education chemistry: emerging trends and potential directions & 2015 & OA & kvalificerar eller motsäger påståendet (0.94) \\
Flipped Learning in Saudi Higher Education during Pandemics & 2022 & OA & kvalificerar eller motsäger påståendet (0.82) \\
Flipped Learning, Flipped Satisfaction: Getting the Balance Right & 2015 & OA & kvalificerar eller motsäger påståendet (0.94) \\
Flipping General Chemistry in Small Classes: Students' Perception and Success & 2019 & OA & kvalificerar eller motsäger påståendet (0.97) (refutes-claim) \\
Flipping Laboratory Sessions in a Computer Science Course: An Experience Report & 2021 & IEEE & kvalificerar eller motsäger påståendet (0.91) \\
Four Pedagogical Dimensions for Understanding Flipped Classroom Practices in Higher Education: A Systematic Review & 2019 & OA & kvalificerar eller motsäger påståendet (0.97) \\
From Theory to Practice: A Systematic Review of the Flipped Classroom in Higher Music Education (2021-2025) & 2025 & Scopus & kvalificerar eller motsäger påståendet (0.78) \\
Gamified Flipped Classroom and Effective Learning Outcomes Through Bibliometric and Systematic Literature Review Toward SDGs 4 & 2026 & OA & kvalificerar eller motsäger påståendet (0.83) \\
Gamified flipped classroom in higher education: a meta-analysis of (moderated) cognitive and non-cognitive effects & 2026 & WoS & kvalificerar eller motsäger påståendet (0.98) \\
Higher education dominance and siloed knowledge: a systematic review of flipped classroom research & 2018 & OA & kvalificerar eller motsäger påståendet (0.91) \\
How Can We Support Students' Learning Experiences in Higher Education? Campus Wide Course Transformation Program Systematic Review and Meta-Analysis & 2022 & WoS & kvalificerar eller motsäger påståendet (0.88) \\
Impact of a hybrid flipped classroom and case-based learning model on learning outcomes and competency development in Epidemiology & 2026 & Scopus & kvalificerar eller motsäger påståendet (0.96) \\
Impact of Active Learning on Object-Oriented Programming Instruction : Transforming from 3D to Text-based coding & 2021 & IEEE & kvalificerar eller motsäger påståendet (0.79) \\
Impact of different teaching modes on medical students' performance under the scoring criteria for multiple-choice questions A meta-analysis & 2024 & Scopus & kvalificerar eller motsäger påståendet (0.78) \\
Impact of flipped classrooms on higher education student outcomes: A second-order meta-analysis & 2026 & Scopus & kvalificerar eller motsäger påståendet (0.99) \\
Impact of integrating paper-based pulmonary function test case group discussions into flipped classroom on residents' COPD grading assessment competency & 2026 & WoS & kvalificerar eller motsäger påståendet (0.94) \\
Impact of the Flipped Classroom on Medical Education: A Systematic Review with Meta-Analysis & 2026 & Scopus & kvalificerar eller motsäger påståendet (0.96) \\
Implementation of a Mixed Strategy of Gamification and Flipped Learning in Undergraduate Basic Programming Courses & 2023 & OA & kvalificerar eller motsäger påståendet (1.00) \\
Implementing peer-assisted learning in a dental flipped classroom: a model for postgraduate-undergraduate collaboration & 2026 & Scopus & kvalificerar eller motsäger påståendet (0.95) \\
Insights from an Umbrella Review of Flipped Learning in Higher Education & 2025 & OA & kvalificerar eller motsäger påståendet (0.97) \\
Introducing an Inverted Classroom Concept for an Undergraduate Course in Power Electronics -- Implementation, Evaluation and Experiences & 2019 & IEEE & kvalificerar eller motsäger påståendet (0.97) (refutes-claim) \\
Inversed Learning in an Intermediate Accounting Course. & 2019 & OA & kvalificerar eller motsäger påståendet (0.86) \\
Inverted Classroom by Topic - A Study in Mathematics for Electrical Engineering Students & 2014 & OA & kvalificerar eller motsäger påståendet (0.98) (refutes-claim) \\
Investigating flipped learning: Post-secondary student selfregulated learning, perceptions, and achievement & 2015 & OA & kvalificerar eller motsäger påståendet (0.91) \\
Investigating Flipped Learning: Student Self-Regulated Learning, Perceptions, and Achievement in an Introductory Biology Course & 2017 & OA & kvalificerar eller motsäger påståendet (0.91) \\
Investigating the impact of the flipped classroom model on higher order thinking skills: A metaanalysis study & 2024 & Scopus & kvalificerar eller motsäger påståendet (0.68) \\
Learner Preparedness in Flipped Classroom: A Case Study of a Flipped Postgraduate Course & 2020 & IEEE & kvalificerar eller motsäger påståendet (0.84) \\
Learning Achievements and Attitudes in a Computer Science Course: Activating Students Flipped Learning via ICT Technologies & 2017 & IEEE & kvalificerar eller motsäger påståendet (0.83) \\
Lectures Versus Flipped-Classroom Learning in Anatomy: Cross-Testing Evidence on Performance, Transfer, and Student Satisfaction & 2025 & WoS & kvalificerar eller motsäger påståendet (0.99) \\
Leveraging active learning through student-generated questions and electronic posters in a physiology course & 2026 & Scopus & kvalificerar eller motsäger påståendet (0.93) \\
Linking Motivation and Learning Strategy Changes to Learning Outcomes in a Flipped Engineering Laboratory & 2026 & Scopus & kvalificerar eller motsäger påståendet (0.67) \\
LOOKING INTO THE IMPACT OF FLIPPED LEARNING PEDAGOGY ON STUDENTS' PERCEIVED LEARNING EXPERIENCE IN UNDERGRADUATE MATHEMATICS COURSES & 2017 & WoS & kvalificerar eller motsäger påståendet (0.97) \\
Merging Studio-Based Learning With Flipped Classroom Techniques: Ensuring Support and Focus in Engineering Education & 2025 & IEEE & kvalificerar eller motsäger påståendet (0.91) \\
Meta-Analisis investigasi model kelas terbalik pada keterampilan berpikir tingkat tinggi (HOTS) siswa matematika: Analisis efek gabungan dan heterogenitas & 2024 & OA & kvalificerar eller motsäger påståendet (0.92) \\
Meta-Analysis of Flipped Classroom on Students' Mathematics Abilities: Effectiveness and Heterogeneity Analysis & 2023 & OA & kvalificerar eller motsäger påståendet (0.91) \\
Meta-analytic Review of Flipped Classrooms in Health Professions Education & 2026 & Scopus & kvalificerar eller motsäger påståendet (1.00) \\
Methodological gaps and domain imbalances in flipped learning reviews: A critical umbrella review & 2026 & WoS & kvalificerar eller motsäger påståendet (0.99) \\
Mixed method evaluation of a novel seminar format for a PBL-based integrated curriculum & 2026 & WoS & kvalificerar eller motsäger påståendet (0.84) \\
MOOC-Based Flipped Classroom for On-Campus Teaching in Undergraduate Engineering Courses & 2023 & WoS & kvalificerar eller motsäger påståendet (0.91) \\
MOTIVATION IN ENGLISH LANGUAGE TEACHING AND LEARNING THROUGH THE ICT-MEDIATED FLIPPED CLASSROOM MODEL: A SYSTEMATIC LITERATURE REVIEW IN INTERNATIONAL CONTEXTS & 2025 & Scopus & kvalificerar eller motsäger påståendet (0.82) \\
Moving from Passive Learning to Active Learning: Investigating the Effect of the Flipped Classroom Pedagogical Approach in a Computer-Based Technology Course in Higher Education & 2025 & Scopus & kvalificerar eller motsäger påståendet (0.72) \\
Módszertani keretrendszerek a flipped classroom kutatásaiban & 2025 & OA & kvalificerar eller motsäger påståendet (0.90) \\
Módszertani keretrendszerek a flipped classroom kutatásaiban = Methodological frameworks in flipped classroom research & 2025 & OA & kvalificerar eller motsäger påståendet (0.98) \\
Nursing students' outcomes in the flipped classroom approach: An integrative review & 2025 & Scopus & kvalificerar eller motsäger påståendet (0.96) \\
Paradoxical effect of flipped classroom on nursing students' learning ability and satisfaction in a fundamental of nursing clinical course: A quasi-experimental study & 2025 & Scopus & kvalificerar eller motsäger påståendet (0.98) \\
PERCEPTION OF FLIPPED LEARNING IN MONGOLIAN HIGHER EDUCATION: A QUANTITATIVE STUDY &  & OA & kvalificerar eller motsäger påståendet (0.86) \\
Pharmacy students' perception of learning and engagement in a flipped-classroom of a physiology course & 2021 & OA & kvalificerar eller motsäger påståendet (0.91) \\
Preliminary Assessment of the Effect of a Flipped Classroom Combined with Mind Mapping on Learning Outcomes in Ultrasound Use in Animal Husbandry Education: A Comparison with Traditional Lecture-Based Learning Among Third-Year BS Students in China & 2026 & Scopus & kvalificerar eller motsäger påståendet (0.95) \\
Probing the Inverted Classroom: A Controlled Study of Teaching and Learning Outcomes in Undergraduate Engineering and Mathematics & 2020 & OA & kvalificerar eller motsäger påståendet (0.87) (refutes-claim) \\
Putting Flipped Classroom into Practice: A Comprehensive Review of Empirical Research & 2018 & Scopus & kvalificerar eller motsäger påståendet (0.83) \\
Regulation of Flipped Learning Activities in Programming & 2021 & OA & kvalificerar eller motsäger påståendet (0.90) \\
Reinforcing General Chemistry Learning with Structured Daily Recaps: A Pilot Study & 2026 & Scopus & kvalificerar eller motsäger påståendet (0.94) \\
Review of Studies on the Use of the Flipped Learning Model in Science Education & 2025 & WoS & kvalificerar eller motsäger påståendet (0.81) \\
Revisiting flipped class model effectiveness on learning performance: A meta-analysis of meta analyses & 2023 & OA & kvalificerar eller motsäger påståendet (0.99) \\
Seeking the best blend for deep learning in a flipped classroom -- viewing student perceptions through the Community of Inquiry lens & 2018 & OA & kvalificerar eller motsäger påståendet (0.87) \\
Self-Efficacy, Emotions and Mathematics Achievement: Longitudinal Effects of Flipped Learning & 2025 & Scopus & kvalificerar eller motsäger påståendet (0.91) \\
South African students' perceptions of the flipped classroom: A case study of higher education & 2020 & OA & kvalificerar eller motsäger påståendet (0.96) \\
Student Perceptions of a Flipped Pharmacotherapy Course & 2015 & OA & kvalificerar eller motsäger påståendet (0.96) \\
Student Perceptions of Active Learning. & 2015 & OA & kvalificerar eller motsäger påståendet (0.82) \\
Student Perceptions of Flipped Learning & 2015 & OA & kvalificerar eller motsäger påståendet (0.91) \\
Student Perceptions of Flipped Learning in STEM Courses & 2026 & WoS & kvalificerar eller motsäger påståendet (0.79) \\
Student Perceptions of Inverted Classroom Benefits in a First-Year Engineering Course & 2020 & OA & kvalificerar eller motsäger påståendet (0.83) \\
Student Perceptions of Lecture-Capture Video to Facilitate Learning in a Flipped Classroom & 2018 & OA & kvalificerar eller motsäger påståendet (0.89) \\
Student Perceptions of the Flipped Classroom Using Edpuzzle and the Zoom Meeting App in Undergraduate Mathematical Physics Course & 2022 & IEEE & kvalificerar eller motsäger påståendet (0.86) \\
Student Perceptions of the Impact of Using the Flipped Classroom Approach for an Introductory-level Multidisciplinary Module & 2014 & OA & kvalificerar eller motsäger påståendet (0.76) \\
Student Satisfaction: Online Learning-based MIKiR Approach in UNRI-UIN SUSKA RIAU & 2021 & IEEE & kvalificerar eller motsäger påståendet (0.72) \\
Student-Perceived Effectiveness of Online Content Delivery Modes & 2015 & WoS, Scopus & kvalificerar eller motsäger påståendet (0.78) \\
Students' behavioural, cognitive and affective outcomes in gamified flipped classrooms: A meta-analysis & 2025 & Scopus & kvalificerar eller motsäger påståendet (0.94) \\
Students' experiences of active engagement through cooperative learning activities in lectures & 2011 & OA & kvalificerar eller motsäger påståendet (0.80) \\
Students' perception of flipped classroom: A case study for private universities in Jordan & 2019 & OA & kvalificerar eller motsäger påståendet (0.88) \\
Students' Perception of Flipped Classroomin in Engineering Education & 2024 & IEEE & kvalificerar eller motsäger påståendet (0.85) \\
Students' perceptions of using a flipped classroom instructional model in an ESP course & 2017 & IEEE & kvalificerar eller motsäger påståendet (0.88) \\
Switching to blend-Ed: Effects of replacing the textbook with the browser in an introductory computer programming course & 2016 & IEEE & kvalificerar eller motsäger påståendet (0.96) (refutes-claim) \\
Systematic Literature Review of Flipped Classroom in Mathematics & 2021 & OA & kvalificerar eller motsäger påståendet (0.99) \\
Systematic Review and Results of the Experiment of a Flipped Learning Model for the Courses of Descriptive Geometry, Engineering and Computer Graphics, Computer Geometry & 2017 & OA & kvalificerar eller motsäger påståendet (0.87) \\
Teaching Electromagnetism with the Inverted Classroom Approach: Student Perceptions and Lessons Learned & 2020 & OA & kvalificerar eller motsäger påståendet (0.83) \\
Teaching Programming in the Age of Generative Artificial Intelligence: Learning Gains and Pedagogical Integration in a Higher Education Context & 2026 & Scopus & kvalificerar eller motsäger påståendet (0.92) \\
Teaching Tip: The Flipped Classroom. & 2014 & OA & kvalificerar eller motsäger påståendet (0.70) \\
Team-based learning versus lecture-based instruction for chest radiograph interpretation in physician associate education: a quasi-experimental study & 2026 & Scopus & kvalificerar eller motsäger påståendet (0.80) \\
Technology-enhanced differentiation in public administration education: A flipped classroom approach to student heterogeneity & 2026 & WoS & kvalificerar eller motsäger påståendet (0.88) \\
The Effect and Strategies of Flipped Learning in Nursing Education: A Systematic Review & 2018 & OA & kvalificerar eller motsäger påståendet (0.72) \\
The Effect Of Flipped Classroom Model On Students' Academic Achievement In Science And Mathematics Education: A Meta-Analysis Study & 2021 & OA & kvalificerar eller motsäger påståendet (0.99) \\
The Effect of Motivation on Learners' Performance and Satisfaction under Flipped Strategy in Discrete Math & 2020 & IEEE & kvalificerar eller motsäger påståendet (0.97) (refutes-claim) \\
The Effectiveness of Flipped Classroom in English Language Learning: A Meta-Analysis & 2025 & Scopus & kvalificerar eller motsäger påståendet (0.98) \\
The Effectiveness of the Flipped Classroom Method on Students' Speaking Ability: A Meta-Analysis Study & 2023 & OA & kvalificerar eller motsäger påståendet (0.90) \\
The effectiveness of the flipped classroom: A second-order meta-analysis & 2025 & Scopus & kvalificerar eller motsäger påståendet (0.99) \\
The Effectiveness of Using MOOC with E-FLIC Approach in Remote Learning as New Normal of Education to Support Basic Computer Science Learning & 2025 & IEEE & kvalificerar eller motsäger påståendet (0.98) (refutes-claim) \\
The effects of classroom flip on the student learning experience: An investigative study in UAE classrooms & 2013 & IEEE & kvalificerar eller motsäger påståendet (0.91) \\
The effects of flipped classrooms to improve learning outcomes in undergraduate health professional education: A systematic review & 2023 & OA & kvalificerar eller motsäger påståendet (1.00) \\
The evidence for 'flipping out': A systematic review of the flipped classroom in nursing education & 2016 & WoS & kvalificerar eller motsäger påståendet (0.96) \\
The Flipped Classroom & 2013 & OA & kvalificerar eller motsäger påståendet (0.88) \\
The Flipped Classroom in Medical Education: Systematic Review and Meta-Analysis & 2025 & OA & kvalificerar eller motsäger påståendet (0.94) \\
The flipped classroom in second language learning: A meta-analysis & 2023 & Scopus & kvalificerar eller motsäger påståendet (0.99) \\
The flipped classroom: A learning model to increase student engagement not academic achievement & 2017 & OA & kvalificerar eller motsäger påståendet (0.99) (refutes-claim) \\
The Flipped Classroom: A Survey of the Research & 2020 & OA & kvalificerar eller motsäger påståendet (0.86) \\
The flipped classroom: supporting a diverse group of students in their learning & 2019 & OA & kvalificerar eller motsäger påståendet (0.91) \\
The impact of active methodologies in the classroom on academic achievement: A systematic literature review using the PRISMA protocol & 2025 & WoS & kvalificerar eller motsäger påståendet (0.87) \\
The Impact of Flipped Classroom in Undergraduate Nursing Education & 2025 & OA & kvalificerar eller motsäger påståendet (0.88) \\
The Impact of Flipped Learning on L2 Learners' Achievements: A Meta-Analysis & 2023 & OA & kvalificerar eller motsäger påståendet (0.96) \\
The Impact of Gamification on Nursing Students' Academic Performance: An Integrative Review & 2026 & Scopus & kvalificerar eller motsäger påståendet (0.88) \\
The impact of the flipped classroom in a principles of microeconomics course: evidence from a quasi-experiment with two flipped classroom designs & 2018 & OA & kvalificerar eller motsäger påståendet (0.99) \\
The inverted classroom : a literature review & 2012 & OA & kvalificerar eller motsäger påståendet (0.88) \\
The relationship between perceived learning and student attainment: lectures and problem-based learning sessions & 2017 & OA & kvalificerar eller motsäger påståendet (0.90) \\
The right tool for the job: matching active learning techniques to learning objectives & 2024 & OA & kvalificerar eller motsäger påståendet (0.87) \\
The use of voiceovers and review games in teaching immunology & 2026 & Scopus & kvalificerar eller motsäger påståendet (0.97) (refutes-claim) \\
Toward a set of design principles for mathematics flipped classrooms: A synthesis of research in mathematics education & 2017 & OA & kvalificerar eller motsäger påståendet (0.98) \\
Traditional lectures versus flipped classroom in manual therapy education: a cross-sectional study of student engagement and learning perceptions & 2026 & WoS & kvalificerar eller motsäger påståendet (0.96) \\
Transforming Undergraduate Chemistry: A Decade of Flipped Classroom Innovations (2013-2023) & 2025 & Scopus & kvalificerar eller motsäger påståendet (0.91) \\
Trends of flipped classroom studies for physics learning: A systematic review & 2021 & Scopus & kvalificerar eller motsäger påståendet (0.95) \\
Undergraduate medical students' perceptions of an interactive and collaborative cloud-based learning strategy: survey at a single institution & 2026 & Scopus & kvalificerar eller motsäger påståendet (0.91) \\
Undergraduates' satisfaction and perceptions of learning outcomes across teacher- and learner-focused pedagogies & 2019 & Scopus & kvalificerar eller motsäger påståendet (0.97) (refutes-claim) \\
Unfolding the Student Experience: Analysis of Flipped Classrooms Post-COVID-19 & 2026 & Scopus & kvalificerar eller motsäger påståendet (0.58) \\
University Students' Perception of the Usefulness of the Flipped Classroom Methodology & 2020 & OA & kvalificerar eller motsäger påståendet (0.89) \\
Use of Gamification as an Active Teaching and Learning Methodology in Undergraduate Dental Education in Oral Radiology & 2026 & WoS & kvalificerar eller motsäger påståendet (0.86) \\
Using Flipped Classroom and Team-Based Learning in a First-Semester Programming Course: An Experience Report & 2019 & IEEE & kvalificerar eller motsäger påståendet (0.97) \\
Using flipped classroom approach to teach computer programming & 2016 & IEEE & kvalificerar eller motsäger påståendet (0.87) \\
Using Modern Learning Method to Teach Pharmacy Students Psychopharmacotherapy & 2024 & WoS & kvalificerar eller motsäger påståendet (0.80) \\
What makes flipped learning effective? A meta-analysis of teaching and learning factors & 2026 & Scopus & kvalificerar eller motsäger påståendet (0.96) \\
What millennial medical students say about flipped learning & 2017 & OA & kvalificerar eller motsäger påståendet (0.98) \\
When is flipped classroom more effective and why it flops & 2020 & Scopus & kvalificerar eller motsäger påståendet (0.77) \\
Whether case-based teaching combined with the flipped classroom is more valuable than traditional lecture-based teaching methods in clinical medical education: a systematic review and meta-analysis & 2025 & Scopus & kvalificerar eller motsäger påståendet (0.95) \\
\end{longtable}
\endgroup

